\documentclass{article}
\usepackage{iclr2027_conference,times}
\input{math_commands.tex}
\usepackage{amsmath,amssymb,booktabs,microtype,multirow}
\usepackage{hyperref}
\usepackage{url}

\title{Characteristic Spectral Transform Distillation:\\
Projector-Free Multi-Scale Geometry Matching for Language Models}

\author{Anonymous Authors}

% Keep commented for double-blind submission.
% \iclrfinalcopy

\newcommand{\cst}{\textsc{CST}}
\newcommand{\sg}{\operatorname{sg}}

\begin{document}
\maketitle

\begin{abstract}
Knowledge distillation for language models usually constrains output
distributions while leaving intermediate representation geometry
underspecified.  Existing geometry-aware objectives often require learnable
width-matching projectors, stored teacher prototypes, or explicit spectral
decompositions.  We introduce \emph{Characteristic Spectral Transform}
(\cst) distillation, a projector-free objective that compares student and
teacher hidden states through the multi-scale transform
$\Phi_H(\gamma)=\log\det(I+\gamma G_H)$, where $G_H$ is a centered,
trace-normalized Gram matrix.  Sampling several positive scales $\gamma$
makes the loss sensitive to the complete normalized spectrum without directly
matching eigenvalues or requiring equal hidden widths.  A practical estimator
uses a shared response-token subset, reuses one Gram matrix across scales, and
evaluates all scales with batched Cholesky factorizations.  This yields an
auxiliary objective with no inference-time parameters and negligible measured
kernel overhead.  In a preliminary 10K-example distillation of
Qwen2.5-14B-Instruct into Qwen2.5-1.5B-Instruct, \cst{} improves the
eight-task reasoning average from 44.6 for the reference base student to 49.9
under legacy-compatible metrics, while using only a fraction of the reference
training budget.  We also identify evaluation-version sensitivity on MATH and
report both exact-string and symbolic-verification scores.  These results
motivate transform-level spectral supervision as a simple and scalable
alternative to projector-based hidden-state alignment.
\end{abstract}

\section{Introduction}

Autoregressive knowledge distillation (KD) transfers behavior from a large
teacher to a smaller student by aligning their predictive distributions
\citep{hinton2015distilling,kim2016sequence,gu2024minillm,agarwal2024gkd}.
Modern on-policy methods improve this transfer by training on student
trajectories, reducing the mismatch between training and generation
\citep{ko2024distillm}.  Nevertheless, an output objective observes hidden
representations only through the final vocabulary head.  Many different
intermediate geometries can therefore induce similar token distributions.

Matching hidden states directly is not an entirely satisfactory remedy.
Teacher and student widths differ, making pointwise losses dependent on a
learned projector whose parameters and inductive biases become part of the
objective.  Distributional summaries based on centroids add a precomputation
stage and state that must be stored during training.  Explicit singular-value
or eigenvalue matching exposes the training path to expensive decompositions
and unstable gradients around repeated eigenvalues.  More fundamentally,
matching a finite ordered spectrum treats the eigenvalues themselves as the
representation, rather than comparing a characteristic functional of the
geometry.

We propose \cst{} distillation.  For response-token hidden states
$H\in\R^{m\times d}$, we center across tokens, normalize total energy, and
form a positive semidefinite Gram matrix $G_H$ with unit trace.  We then
evaluate
\begin{equation}
  \Phi_H(\gamma)=\log\det(I+\gamma G_H),\qquad \gamma>0,
  \label{eq:transform}
\end{equation}
at multiple randomly sampled scales.  The student learns by matching the
teacher's transform values at the same scales and token positions.  The
construction is invariant to translation, isotropic scaling, and orthogonal
changes of hidden coordinates, and it naturally compares representations of
different widths.

Our contributions are:
\begin{itemize}
  \item We formulate hidden-state distillation as matching a multi-scale
  log-determinant transform of normalized geometry.  The transform identifies
  the finite normalized spectrum over any nontrivial interval, but the
  training loss never computes an eigendecomposition.
  \item We give a width-agnostic, projector-free estimator using shared token
  and scale samples, the smaller determinant space, Gram reuse, and batched
  Cholesky factorization.
  \item We integrate \cst{} into an on-policy output-KD pipeline and provide
  numerical, invariance, gradient, degeneracy, and dispatch tests.  Profiling
  on an RTX PRO 6000 Blackwell reports approximately 2.2ms for the \cst{}
  kernel at four layers, 64 tokens, and two scales.
  \item We report a transparent preliminary Qwen study, including the
  substantial discrepancy between legacy exact-string and modern symbolic
  MATH evaluation, and state the additional controlled experiments required
  for a definitive comparison.
\end{itemize}

\section{Related Work}

\paragraph{Language-model distillation.}
Classical KD minimizes a softened output divergence
\citep{hinton2015distilling}, while sequence-level KD trains on teacher
sequences \citep{kim2016sequence}.  On-policy approaches instead expose the
teacher to student-generated trajectories.  MiniLLM uses reverse KL
\citep{gu2024minillm}; GKD generalizes the divergence family
\citep{agarwal2024gkd}; and DistiLLM combines adaptive sampling and skewed
divergences \citep{ko2024distillm}.  \cst{} is complementary to these methods:
it leaves output KD unchanged and adds intermediate geometric supervision.

\paragraph{Representation transfer.}
Intermediate-feature distillation has long used layer mappings and learned
projections to reconcile dimensionality \citep{romero2015fitnets,sun2019patient}.
Relational objectives transfer pairwise or higher-order structure instead of
coordinates \citep{park2019relational,tung2019similarity}.  Recent spectral
approaches diagnose or regularize representation collapse through rank,
nuclear norm, or entropy summaries \citep{roy2007effective,jing2022dimensional}.
Our distinction is operational and conceptual: \cst{} matches a
characteristic transform, uses neither a learnable width adapter nor a teacher
prototype bank, and avoids explicit spectrum computation in training.

\paragraph{Log determinants.}
Log-determinant functionals occur in Gaussian information measures,
experimental design, and diversity objectives.  If $p_i$ are the eigenvalues
of a trace-normalized PSD matrix, Equation~\ref{eq:transform} equals
$\sum_i\log(1+\gamma p_i)$.  Rather than maximizing this quantity at a fixed
scale, we use its curve over scale as a signature that the student matches to
the teacher.

\section{Characteristic Spectral Transform Distillation}

\subsection{Normalized hidden-state geometry}

Let $H\in\R^{m\times d}$ contain the hidden states of $m$ selected response
tokens at one layer.  Define
\begin{align}
  X &= H-\mathbf{1}\bar h^\top, &
  \bar h &= \frac{1}{m}\sum_{i=1}^m H_{i,:},\\
  e_H &= \lVert X\rVert_F^2, &
  G_H &= \frac{XX^\top}{e_H+\epsilon}.
  \label{eq:gram}
\end{align}
For nondegenerate $X$, $G_H\succeq0$ and $\Tr(G_H)\approx1$.  Centering
removes global translation, while energy normalization removes isotropic
scale.  Because $XX^\top$ is unchanged by $X\mapsto XQ$ for orthogonal $Q$,
the construction does not require teacher and student coordinates to align.

\subsection{A multi-scale spectral signature}

For $gamma>0$, \cst{} is Equation~\ref{eq:transform}.  Writing the positive
eigenvalues of $G_H$ as $p_1,\ldots,p_r$ gives the analysis identity
\begin{equation}
  \Phi_H(\gamma)=\sum_{i=1}^{r}\log(1+\gamma p_i).
  \label{eq:spectral-identity}
\end{equation}
Equation~\ref{eq:spectral-identity} explains the signal but is \emph{not} the
training algorithm.  Small scales emphasize low-order spectral moments,
whereas large scales increasingly expose rank and the distribution of energy
among active directions.  In particular,
\begin{equation}
 \Phi_H(\gamma)
 =\gamma-\frac{\gamma^2}{2}\Tr(G_H^2)+O(\gamma^3),
 \label{eq:taylor}
\end{equation}
because $\Tr(G_H)=1$.  The first nonconstant discriminative term therefore
measures spectral concentration.

The transform is characteristic for a finite spectrum.  Differentiating,
\begin{equation}
 \Phi'_H(\gamma)=\sum_i\frac{p_i}{1+\gamma p_i},
\end{equation}
a rational function whose poles and residues determine the nonzero $p_i$.
Thus equality of transforms on an open positive interval implies equality of
the normalized nonzero spectra, including multiplicities.  Finite Monte Carlo
sampling provides a stochastic approximation to this population criterion.

\subsection{Distillation objective}

For matched student and teacher layers $\ell_s$ and $\ell_t$, we sample the
same valid response-token indices and scales
\begin{equation}
 \log\gamma_j\sim
 \mathcal U(\log\gamma_{\min},\log\gamma_{\max}),
 \quad j=1,\ldots,q.
\end{equation}
The per-layer loss is
\begin{equation}
 \mathcal L_{\cst}^{(\ell)}=
 \frac{1}{q}\sum_{j=1}^{q}
 \left[
   \Phi_{H_s^{(\ell_s)}}(\gamma_j)
   -\sg\!\left(\Phi_{H_t^{(\ell_t)}}(\gamma_j)\right)
 \right]^2,
\end{equation}
and we average uniformly across $n$ selected layers.  The complete objective is
\begin{equation}
 \mathcal L=\mathcal L_{\mathrm{output\text{-}KD}}
 +\lambda_{\cst}\mathcal L_{\cst}.
 \label{eq:total}
\end{equation}
No student-to-teacher hidden projector appears in this path.  The teacher
transform is computed under stop-gradient, and \cst{} introduces no
inference-time module.

\subsection{Efficient and stable computation}

We do not call a determinant, inverse, SVD, or eigendecomposition.  Given the
centered $X$, we construct only the smaller Gram space.  When $m\leq d$ we use
$B=XX^\top/(e_H+\epsilon)$.  If $d<m$, Sylvester's determinant theorem gives
\begin{equation}
 \det(I_m+\gamma XX^\top/e_H)
 =\det(I_d+\gamma X^\top X/e_H),
\end{equation}
so we instead use $B=X^\top X/(e_H+\epsilon)$.  After symmetrization, all
$q$ matrices are constructed at once:
\begin{equation}
 A_j=I+\gamma_jB,\qquad
 A\in\R^{q\times r\times r},\quad r=\min(m,d).
\end{equation}
A single batched Cholesky call produces $A_j=L_jL_j^\top$, followed by
\begin{equation}
 \Phi_H(\gamma_j)=2\sum_k\log (L_j)_{kk}.
\end{equation}
All spectral arithmetic is float32.  We retry with a small diagonal jitter
only if Cholesky reports failure and skip layers with negligible centered
energy.  The dominant per-layer costs are $O(md r)$ for the smaller Gram and
$O(qr^3)$ for Cholesky, with $r=m=64$ in our default LLM setting.

\begin{table}[t]
\caption{Default \cst{} training configuration.  One token subset and one set
of scales are shared by teacher, student, and all supervised layers.}
\label{tab:config}
\centering
\small
\begin{tabular}{lc@{\qquad}lc}
\toprule
Maximum response tokens $m$ & 64 & Supervised layers $n$ & 4\\
Gamma samples $q$ & 2 & Layer-depth range & $[0.20,0.85]$\\
$\gamma_{\min},\gamma_{\max}$ & $10^{-2},10^2$ & Sampling & log-uniform\\
Distance & squared $\ell_2$ & Arithmetic & float32\\
\bottomrule
\end{tabular}
\end{table}

\section{Experimental Setup}

\paragraph{Models and data.}
Our preliminary experiment distills Qwen2.5-14B-Instruct into
Qwen2.5-1.5B-Instruct \citep{qwen2024qwen25}.  Training uses 10K prompts from
UltraInteract \citep{yuan2024ultrainteract}, rather than the full roughly 78K
budget used by the reference study.  Student responses are generated on
policy.  We train for two epochs with global batch size 64, AdamW, learning
rate $5\times10^{-6}$, cosine decay, 10\% warmup, maximum length 1024, and
output-KD temperature 1.  Four uniformly spaced layers in relative depth
$[0.20,0.85]$ are supervised.  Based on loss-scale calibration, we use
$\lambda_{\cst}=0.003$; the raw transform loss has a different scale from
nuclear-norm losses and should not inherit their coefficient.

\paragraph{Evaluation.}
We use lm-evaluation-harness with greedy decoding and the model chat template.
The suite contains GSM8K, GSM-Plus, MATH, MBPP, SciQ, MMLU-STEM,
MMLU-Pro-Math, and three-shot BBH-CoT.  Evaluation runs in bfloat16 through
vLLM on one RTX PRO 6000 Blackwell GPU.  For compatibility with the published
table we report the harness's exact-match MATH score in the main average.  The
current harness also provides symbolic \texttt{math\_verify}; because it
labels many mathematically equivalent answers that exact matching rejects, we
report it separately and do not substitute it into the legacy average.

\paragraph{Baselines and reporting status.}
Student, CSD, and NuNo-KD numbers from the reference NuNo-KD manuscript
\citep{anonymous2026nuno} are included only as contextual
references, not as controlled reproductions.  Their training budget and, for
MATH, evaluator version differ from ours.  A final submission must replace
these comparisons with checkpoints trained on the same data and evaluated by
the same frozen harness.  We include this explicit limitation to prevent the
preliminary results from being interpreted as a definitive ranking.

\section{Results}

\begin{table*}[t]
\caption{Preliminary Qwen2.5-14B $\rightarrow$ 1.5B reasoning accuracy (\%).
Reference rows are contextual and use the NuNo-KD manuscript's full training
budget.  Our \cst{} row uses 10K examples.  MATH uses legacy-compatible exact
match; \texttt{math\_verify} for the same \cst{} generations is 45.06.}
\label{tab:main-results}
\centering
\small
\setlength{\tabcolsep}{3.0pt}
\begin{tabular}{lrrrrrrrrr}
\toprule
Method & GSM8K & GSM+ & MATH & MBPP & SciQ & STEM & Pro-Math & BBH & Avg.\\
\midrule
Student (reference) & 57.1 & 38.8 & 16.9 & 38.4 & 85.5 & 49.5 & 34.4 & 36.5 & 44.6\\
CSD (reference) & 61.5 & 41.5 & 29.3 & 43.8 & 81.9 & 51.3 & 44.4 & 42.1 & 49.5\\
NuNo-KD (reference) & 65.2 & 46.8 & 32.2 & 44.4 & 86.2 & 51.7 & 42.0 & 44.9 & 51.7\\
\midrule
\cst{} (ours, 10K) & \textbf{65.58} & 46.79 & 14.92 & \textbf{46.80} & \textbf{86.90} &
\textbf{51.95} & \textbf{44.49} & 41.87 & 49.91\\
\bottomrule
\end{tabular}
\end{table*}

Despite using 10K rather than roughly 78K prompts, \cst{} improves the
reference base-student average by 5.3 points and reaches 49.91 under the
legacy-compatible metric selection (Table~\ref{tab:main-results}).  Its
GSM8K score is 65.58, with standard error 1.31 points, statistically
indistinguishable from the reference NuNo-KD value of 65.2.  Gains are also
visible on MBPP, SciQ, MMLU-STEM, and MMLU-Pro-Math.  BBH-CoT remains below the
reference NuNo-KD score.

MATH requires special care.  Current-harness exact match is 14.92, whereas
symbolic verification gives 45.06 on identical generations.  Among 4,454
available per-sample records, 1,409 are rejected by exact match but accepted by
symbolic verification; inspection shows many contain a correct boxed answer
without the exact expected ``Final Answer'' surface form.  Only four examples
show the reverse disagreement.  This establishes strong evaluator sensitivity
but does not establish that 45.06 is comparable to the published 32.2.  The
scientifically valid resolution is to evaluate every checkpoint using both
frozen metrics.

\paragraph{Efficiency.}
At $m=64$, $q=2$, and $n=4$, the reported \cst{} transform computation is
approximately 2.2ms, while an optimizer step in the complete on-policy KD
pipeline takes roughly 22.6s.  This timing isolates the transform kernel and
does not imply that hidden-state capture is free; nevertheless, it indicates
that batched log-determinant evaluation is not the runtime bottleneck.  The
model uses about 31.6GiB allocated GPU memory after initialization in the
reported setup and encounters no \cst{}-related out-of-memory event.

\section{Analysis and Discussion}

\subsection{Why multiple scales?}
At a single small $\gamma$, unit trace makes transforms nearly identical.  At
one large $\gamma$, the signal can be dominated by effective rank and very
small eigenvalues.  Log-uniform sampling over four orders of magnitude mixes
these regimes without constructing a dense grid.  Crucially, teacher and
student receive identical samples, so stochasticity does not introduce a
scale mismatch.  An ablation over $q\in\{1,2,4\}$ and narrower scale intervals
is still needed to quantify this design choice.

\subsection{Relation to rank-profile matching}
Rank-profile matching explicitly eigendecomposes a normalized Gram matrix,
sorts eigenvalues, and compares cumulative spectra.  It is a useful analysis
and baseline tool, but it is not \cst{}.  \cst{} compares evaluations of an
analytic transform and obtains gradients through Cholesky factorization.  It
therefore avoids eigenvector sensitivity at repeated eigenvalues and does not
turn the method into direct eigenvalue regression.

\subsection{Loss-scale calibration}
Raw \cst{} losses in early training are commonly tens of units because the
squared gap aggregates transform differences across broad random scales.  A
coefficient copied from a differently normalized representation objective can
therefore dominate output KD.  In our run, $\lambda_{\cst}=0.2$ would
contribute roughly 12 units when output KD is about 1.4.  We instead select
0.003, yielding a typical contribution of 8--16\% of output loss.  Future work
should tune this value on a held-out validation protocol or normalize the
distance while preserving the transform-level objective.

\section{Limitations and Required Confirmatory Experiments}

The current evidence is preliminary.  First, the 10K training budget differs
from the full reference budget, and published baseline numbers are not
same-run controls.  Second, we report one seed and one model pair.  Third,
MATH metrics changed across harness versions.  Fourth, the present objective
samples response tokens globally across a batch; while indices are shared
across models and layers, alternative per-sequence sampling may change the
geometry.  Finally, transform equality is identifying over a continuum, while
training observes only two random scales per step.

A confirmatory study should include: (i) full-budget Qwen and Gemma2 model
pairs; (ii) output-KD, NuNo-KD, rank-profile, and \cst{} trained with identical
seeds and data; (iii) at least three seeds for the principal comparison;
(iv) ablations over $m\in\{32,64,128\}$, $q\in\{1,2,4\}$,
$n\in\{2,4\}$, gamma range, centering, and loss weight; (v) wall-clock and
peak-memory comparisons including hidden-state capture; and (vi) frozen
evaluation containers with both exact and symbolic MATH scoring.

\section{Conclusion}

We introduced \cst{}, a projector-free representation-distillation objective
based on the characteristic curve $\log\det(I+\gamma G_H)$.  Matching this
curve at shared random scales transfers normalized spectral geometry without
matching coordinates, widths, or eigenvalues directly.  The estimator uses
small response-token subsets and batched Cholesky factorizations, making its
measured kernel cost negligible within on-policy LLM distillation.  A
preliminary 10K-example Qwen experiment is promising but also exposes the need
for controlled full-budget baselines and frozen MATH evaluation.  We view
these transparent intermediate findings as evidence that transform-level
geometry is a useful direction, rather than as a final claim of superiority.

\section*{Ethics Statement}
This work studies model compression and may reduce inference resource use.
Distillation can also transfer teacher errors, biases, and unsafe behaviors;
the proposed representation objective does not mitigate these risks by
itself.  Benchmark accuracy should not be interpreted as evidence of safety or
fitness for deployment.

\section*{Reproducibility Statement}
The implementation exposes token count, layer range, gamma distribution,
number of scale samples, distance, centering, jitter, and loss weight.  Tests
cover Cholesky--log-determinant agreement, Sylvester's identity, invariances,
shared gamma samples, unequal hidden widths, teacher stop-gradient, degenerate
representations, and absence of projector/NuNo state in dispatch.  Exact
commands and evaluator outputs should accompany a submission; the present
draft marks results that still require controlled reproduction.

\section*{AI Use Statement}
Generative AI tools assisted with code review, experiment-log summarization,
and preparation of an initial manuscript draft.  The authors are responsible
for verifying all mathematical statements, citations, experimental results,
and final prose, and for ensuring compliance with the conference policy.

\bibliography{references}
\bibliographystyle{iclr2027_conference}

\appendix
\section{Reference implementation pseudocode}
\begin{verbatim}
indices = sample_response_tokens(labels, max_tokens=m)
gamma = exp(uniform(log(gamma_min), log(gamma_max), [q]))
losses = []
for student_layer, teacher_layer in matched_layers:
    Xs = center(student_hidden[student_layer][indices]).float()
    Xt = center(teacher_hidden[teacher_layer][indices]).float().detach()
    Bs = smaller_gram(Xs) / (squared_frobenius(Xs) + eps)
    Bt = smaller_gram(Xt) / (squared_frobenius(Xt) + eps)
    As = I[None] + gamma[:, None, None] * Bs[None]
    At = I[None] + gamma[:, None, None] * Bt[None]
    phi_s = 2 * log(diag(cholesky(As))).sum(-1)
    phi_t = 2 * log(diag(cholesky(At))).sum(-1)
    losses.append(mean((phi_s - stopgrad(phi_t)) ** 2))
return mean(losses)
\end{verbatim}

\section{Additional numerical safeguards}
Hidden states with centered energy below $\epsilon$ are skipped.  Gram
matrices are explicitly symmetrized to suppress floating-point asymmetry.  We
use \texttt{cholesky\_ex} to inspect factorization status and add
$10^{-6}I$ only on failure.  Teacher hidden states and transforms are detached.
All selected layers share response-token indices and gamma probes within an
optimization step.

\section{Per-category MATH scores}
Current-harness symbolic-verification accuracies are 63.77 (algebra), 37.55
(counting and probability), 38.20 (geometry), 26.69 (intermediate algebra),
35.74 (number theory), 60.39 (prealgebra), and 32.05 (precalculus), for an
aggregate of 45.06.  Corresponding exact-string scores aggregate to 14.92.
These values are reported to make evaluator-version effects auditable.

\end{document}
